\section*{Uwagi bibliograficzne}
\addcontentsline{toc}{section}{Uwagi bibliograficzne}
Przedstawiony tekst jest samodzielnym opracowaniem klasycznego argumentu: grupa wolna, swobodne działanie na części sfery, wchłanianie przeliczalnych wyjątków i przedłużenie promieniowe. Nie zakłada ogólniejszej postaci twierdzenia o równoważności dowolnych ograniczonych zbiorów o niepustym wnętrzu. Udowodniona tutaj wersja to podwojenie kuli domkniętej. Szerszy kontekst i inne warianty rozkładów omawiają:
\begin{itemize}
\item Stefan Banach, Alfred Tarski, \emph{Sur la décomposition des ensembles de points en parties respectivement congruentes}, Fundamenta Mathematicae \textbf{6} (1924), 244--277, \href{https://doi.org/10.4064/fm-6-1-244-277}{DOI: 10.4064/fm-6-1-244-277}.
\item Stan Wagon, \emph{The Banach--Tarski Paradox}, Cambridge University Press, 1985, \href{https://www.cambridge.org/core/books/the-banach-tarski-paradox/EA35091E69D798AE0FAE5A4E8386DD08}{katalog wydawcy}.
\end{itemize}
